\begin{lemma}[Four unused missing colors]
Let \((\varphi,T)\) be a minimum bad pair and let \(T_q\) be the prefix
through \(y_q\), where \(1\le q<|E(T)|\).  Then at least four colors are
missing on \(V(T_q)\) and unused on the colored tree-edge identities
\(e_2,\ldots,e_q\); equivalently,
\[
  |\mathcal F(T_q,e,\varphi)|\ge4.
\]
\end{lemma}
\begin{proof}
Let \(P=T_q\) be the actual prefix object supplied by
\noderef{TreeSequencePrefix}.  Since \(q<|E(T)|\), it is a proper initial
segment, so \noderef{TashTreeProperPrefixElementary} says that \(V(P)\)
is \(\varphi\)-elementary.  We now count incidences directly from the
local data carried by the minimum bad pair.

Unfold \noderef{TashTreeMinimumBadPair} and take its bad-pair conjunct.
By \noderef{TashTreeBadPair}, this conjunct contains a \(k\)-triple; by
\noderef{KTriple}, that triple gives
\[
 k\ge \Delta(G)+1,\qquad
 \text{\(\varphi\) proper},\qquad
 \text{\(e_0\) the unique uncolored exact edge.}
\]
The first-edge and first-attachment fields of \noderef{TreeSequence}
identify \(e_1\) with \(e_0\) and its endpoints with the distinct vertices
\(y_0,y_1\).  In the terminology of \noderef{Incident}, \(e_0\) is
therefore incident with each of them.  The definition
\noderef{IsOneEdgeUncolored} says both that \(e_0\) has value
\(\operatorname{none}\) and that every other exact edge identity has a
colored value.

Fix a vertex \(v\), and consider its colored incident exact edge
identities.  Sending such an identity to its color is injective: two
distinct identities in this set share \(v\), hence are adjacent by
\noderef{AdjacentEdges}, and equality of their colors would contradict
\noderef{IsProperEdgeColoring}.  Its image is exactly
\noderef{PresentColors}: membership in that set is, by definition,
witnessed by an incident edge with the stated color, and every colored
incident edge supplies such a witness.  Thus the number of present colors
equals the number of colored incident identities and is at most the number
of all incident identities counted by \noderef{Degree}.  Consequently
\[
  |\varphi(v)|\le d_G(v)\le\Delta(G),
\]
where the last inequality is the defining maximum property in
\noderef{MaximumDegree}.

At \(v=y_0\) or \(y_1\), the colored incident identities are exactly all
incident identities except \(e_0\): that edge is uncolored, and
\noderef{IsOneEdgeUncolored} colors every other identity.  Removing this
one known incident identity from the \(d_G(v)\) identities therefore
sharpens the bound to
\[
 |\varphi(v)|\le d_G(v)-1\le\Delta(G)-1.
\]
Here \(d_G(v)\ge1\) because \(e_0\) is incident with \(v\), so the
subtraction is legitimate.  Since \noderef{MissingColors} is the
complement of the present-color set in the \(k\) colors, each of
\(y_0,y_1\) misses at least
\[
 k-(\Delta(G)-1)\ge2
\]
colors.  At every other displayed vertex the unsharpened bound gives at
least \(k-\Delta(G)\ge1\) missing color.

The prefix has exactly the \(q+1\) pairwise distinct vertices
\(y_0,\ldots,y_q\), by the injectivity field of
\noderef{TreeSequence}; equivalently these are the elements of
\noderef{TreeSequenceVertexSet}.  Its number of missing-color incidences is
therefore at least
\[
  2+2+(q-1)=q+3.
\]
By \noderef{ColoredVertexSetElementary}, no color occurs in the missing
sets of two distinct prefix vertices.  Hence these incidences have distinct
color coordinates, and the union \noderef{TreeMissingColors} has at least
\(q+3\) elements.

Finally, \noderef{TreeUsedColors} is the set of colors on the exact
identities \(e_2,\ldots,e_q\).  There are \(q-1\) such identities, so this
set has cardinality at most \(q-1\), even if some displayed edges share a
color.  Deleting from a set of at least \(q+3\) colors a set of at most
\(q-1\) colors leaves at least four.  By the set-difference definition in
\noderef{TreeFreeColors}, the resulting set is precisely
\(\mathcal F(P,e_0,\varphi)\), which proves the claim.
\end{proof}
